% =======================================================================
\section{Discussion}

\subsection{Recovered organisational levels}
\label{sec:disc-levels}

The three $H_0$ regimes
(individual, tactical, and team) are the named group-count targets of
Section~\ref{sec:cutoffselection}, read from the pooled group-count
curve on every complete frame of the
ten matches (Figure~\ref{fig:cutoff}). Those regimes remain over a
range of $\delta$ rather than appearing only at the adopted cutoffs.
The metre values are specific to this corpus. Another domain still
requires the interaction lengths to be re-derived.

The tactical distance is the point on that curve where the mean group
count lies nearest five, and where frames with at least four groups
are still common (Figure~\ref{fig:cutoff}c). On the 1~Hz diagnostic
subset, the silhouette score has one interior local maximum, at
$7.0$~m. That distance is a small-group alternative. Once frames that
have collapsed to a single group are omitted, the score rises again
through the two-envelope regime (Supplementary
Figure~\ref{fig:diagnostics}). We record that curve and keep the
group-count cutoff.

\subsection{Two \texorpdfstring{$H_1$}{} regimes}
\label{sec:disc-homology}

Scale-attributable homology
detects loops at two of the three recovered levels. Individual-level
loops are frequent, and tactical-level loops are rarer, with longer
filtration lifetime (Tables~\ref{tab:h1single}--\ref{tab:h1multi};
Figure~\ref{fig:h1bars}). The 150-frame subsample does not resolve
formation or dispersion of loops between samples
(Section~\ref{sec:data}).

On the primary match, every finite $H_1$ bar has a simple-cycle
representative, present at birth by construction
(Figure~\ref{fig:cycles}). Detection is not confined to the adopted
cutoff (Section~\ref{sec:sensitivity}). Tactical $H_1$ is detected on the
band $[6,14]$~m, not only at an isolated optimum.

The pipeline therefore reports three $H_0$ regimes and two $H_1$
regimes. A team-level loop needs four centroids. Four centroids occur in
141 of 436{,}648 frames, and none of those frames has a loop
(Remark~\ref{rem:teamnull}). Recovering
team-level loop structure on the typical envelope would require a
different representation, such as spatial density fields or Delaunay
triangulations. Linkage choice changes tactical $H_1$ counts by nearly
an order of magnitude; that caveat is taken up in
Section~\ref{sec:limitations}.

\subsection{Distinct structure at two levels}
\label{sec:disc-distinct}

The two $H_1$ levels carry distinct rather than redundant
information (Section~\ref{sec:complementarity}): Spearman
$\rho=0.234$ on total persistence, one-sided co-occurrence
(Table~\ref{tab:fisher}), and bottleneck distance on the order of
the tactical mean bar. These statistics do not test whether the
tactical diagram is a coarsening of the individual diagram.

\subsection{Construct validity}
\label{sec:disc-validity}

Total persistence moves with on-ball events
(Section~\ref{sec:eventcorr}). At the individual level, an engagement
is followed by a mean drop of $0.253$~m, and build-up by a mean rise
of $0.700$~m. At the tactical level, both classes move in the same
directions as at the individual level. This is a construct-validity check against measurement
noise. Multiple-testing control and a football reading of the shifts
belong to a companion manuscript.

\subsection{Limitations}
\label{sec:limitations}

Linkage selection is a substantive methodological choice for the
composed map (cluster at $\delta$, then Vietoris--Rips on centroids).
Across a 600-frame, four-match comparison sample, single linkage
detects 106 tactical-level $H_1$ loops. Complete linkage detects 802.
Ward's method detects 822. The difference is almost an order of
magnitude. Point-cloud stability theorems do not bound this map when
single-linkage merges jump; characterising diagram stability under
clustering is an open problem we defer to follow-on work. The gap is
large enough to change which scientific picture the data support, not
merely the precision of one estimate. We retain single linkage because
its definition is the most direct reading of spatial proximity. Its
tendency to chain is a likely source of the gap with complete linkage
and Ward's method; resolving that gap is left open.

A more principled selection criterion would come from the dynamical
system itself, rather than from clustering-quality metrics alone. One
candidate is the linkage under which a governing equation for a
persistence functional admits the sparsest representation, in the
sense of sparse identification of nonlinear dynamics. The functional
would be the landscape $L^2$ norm or the tactical $H_1$
total-persistence sequence. Operationalising this would require a
choice of dynamical state, a defence of the recovered library against
the Akaike information criterion (AIC), the Bayesian information
criterion (BIC), or cross-validation, and a check that the persistence
sequence is lower-dimensional under the selected linkage. We defer this
to the forthcoming full-season work. A population-sized sample of
match sequences is what makes the sparsity comparison meaningful.

A landscape-path stability criterion (minimise variation of the
persistence landscape under small perturbations of $\delta$) remains a
future estimator of the tactical cutoff. It is computable with the
landscape infrastructure of Section~\ref{sec:complementarity}. We
defer it to the persistence-landscape companion work.
Section~\ref{sec:sensitivity} reports the $[6,14]$~m tactical
$H_1$ band as an empirical property around the adopted $11.75$~m.

All data analysed here are broadcast-derived tracking at 10~Hz.
Broadcast tracking typically offers lower spatial precision than
higher-frequency optical systems. That may affect the magnitude of
filtration lifetimes reported. The structural findings (the three
$H_0$ regimes, the $H_1$ presence rates, and the sensitivity profiles)
depend on relative rather than absolute spatial precision. We expect,
but have not tested, that they are robust across tracking technologies.
A direct comparison of filtration lifetimes against optical tracking
data remains for future work.

\subsection{Outlook}

Two extensions are in progress. The ten-match evidence base is being scaled
to a full season of Championship matches. The season will characterise
population-level distributions of $H_0$ and $H_1$ counts, filtration
lifetimes, and landscape norms, and will test whether matches carry a
stable tactical signature at population scale. It is also the sample
on which the linkage criterion above can be evaluated, because a sparse
governing equation for a persistence functional needs a population of
match sequences. Separately, persistence landscape dynamics are being
developed so that each match is a path through landscape space. That
representation is the setting in which the landscape-stability cutoff
can be evaluated, with the match-level landscape, not the per-frame
summary, as the resampling unit.
